\documentclass[12pt]{article}
\usepackage[utf8]{inputenc}
\usepackage[T1]{fontenc}
\usepackage{hyperref}
\title{Soo Language Resources for K-12 Curriculum Development}
\author{Compiled by Claude Sonnet 5 (claude-sonnet-5) for Andr\'as Kornai}
\date{}
\begin{document}
\maketitle

\section{Introduction}
Soo (also So or Tepeth, ISO 639-3 teu) is a Kuliak language of northeastern Uganda, spoken on the slopes of Mount Moroto, Mount Kadam and Mount Napak. Sources describe it as moribund, with only a few dozen elderly speakers. Most of the ethnic group has shifted to Karimojong. The sources reported three dialects: Tepes, Kadam (Katam) and Napak (Yog Toŋi). The search found only encyclopedic and catalogue-type pages.

\section{Primary Educational Resources}
None found.

\section{Linguistic and Reference Materials}
\begin{itemize}
\item Soo language, Wikipedia: \url{https://en.wikipedia.org/wiki/Soo_language}
\item So (Tepeth), Wikipedia: \url{https://en.wikipedia.org/wiki/So_(Tepeth)}
\item Soo numerals page, Max Planck Institute for Evolutionary Anthropology: \url{https://lingweb.eva.mpg.de/channumerals/Soo.htm}
\item Endangered Languages Project entry (bibliography page): \url{https://endangeredlanguages.com/lang/5081/bibliography?hl=de}
\end{itemize}

\section{Religious and Community Resources}
\begin{itemize}
\item Joshua Project profile of the Tepeth in Uganda: \url{https://www.joshuaproject.net/people_groups/15348/UG}
\item Joshua Project language page for Soo: \url{https://joshuaproject.net/index.php/languages/teu}
\end{itemize}
These are demographic profiles and give no information on any Soo-language materials.

\section{Digital and Audio Resources}
None found.

\section{Curriculum Development Context}
With only a few dozen elderly speakers and nearly complete shift to Karimojong, school-based use is not realistic at present. Documentation and revitalisation would have to come first. The sources found are largely derived from Wikipedia and are not independent.
\end{document}